\documentclass{report}
\usepackage[utf8]{inputenc}
\usepackage[spanish]{babel}
\usepackage{amsmath}
\usepackage{amssymb}
\usepackage{amsthm}
\usepackage{float}
\usepackage{tikz}
\usetikzlibrary{shapes}

\theoremstyle{definition}
\newtheorem{definicion}{Definición}[chapter]
\theoremstyle{plain}
\newtheorem{teorema}{Teorema}[chapter]
\newtheorem{lema}{Lema}[chapter]
\theoremstyle{remark}
\newtheorem{ejemplo}{Ejemplo}[chapter]

\begin{document}

\chapter{Preliminares}
\label{cap:preliminares}

En este capítulo se presentan las definiciones y conceptos que serán
usados en el resto del trabajo. Primero, en la
sección~\ref{sec:teoria-lenguajes}, se introducen las nociones
básicas de la teoría de lenguajes: alfabetos, cadenas, lenguajes,
gramáticas y la jerarquía de Chomsky. Después, en la
sección~\ref{sec:automatas}, se presenta el autómata finito. Por último,
en las secciones~\ref{sec:complejidad} y~\ref{sec:sat}, se introduce
la complejidad computacional y se presenta el problema de la
satisfacibilidad booleana con sus variantes polinomiales.

\section{Teoría de Lenguajes}
\label{sec:teoria-lenguajes}

En esta sección se presentan los conceptos de la teoría de lenguajes
que sirven de base al contenido de los capítulos posteriores. La sección
se organiza en tres partes. Primero se introducen los conceptos de
alfabeto, cadena y lenguaje, junto con los problemas de la palabra y
del vacío. Después se definen las gramáticas, un mecanismo que permite
describir lenguajes formales. Para finalizar se presenta la jerarquía
de Chomsky, que clasifica los lenguajes formales según las
restricciones en las reglas de producción de las gramáticas que los
generan.

\subsection{Conceptos básicos}

Los conceptos básicos de la teoría de lenguajes formales son: alfabeto,
cadena y lenguaje. Un \textbf{alfabeto}, denotado como $\Sigma$, es un
conjunto finito y no vacío de símbolos. Una \textbf{cadena} es una
sucesión finita de símbolos del alfabeto. Un \textbf{lenguaje} es un
conjunto de cadenas definido sobre un alfabeto. Por ejemplo, el alfabeto
$\Sigma = \{0, 1\}$ está formado por los símbolos $0$ y $1$; $11$ y $101$
son cadenas sobre el alfabeto $\Sigma$, y un ejemplo de lenguaje es el
conjunto de cadenas formadas por $0$ y $1$ que terminan en $1$.

El símbolo $\varepsilon$ representa la \textbf{cadena vacía}. La
\textbf{longitud} de una cadena $w$, denotada $|w|$, es la cantidad de
símbolos que la componen; en particular, $|\varepsilon| = 0$.

Dadas dos cadenas $w_1$ y $w_2$, la \textbf{concatenación} $w_1 w_2$
se obtiene al escribir los símbolos de $w_1$ seguidos de los símbolos
de $w_2$. Dada una cadena $w$ y un entero $n \ge 0$, por convenio
$w^n$ representa la cadena resultante de concatenar $w$ un total de
$n$ veces, con $w^0 = \varepsilon$.

Si $\Sigma$ es un alfabeto, $\Sigma^+$ denota el conjunto de todas las
cadenas de uno o más símbolos que se pueden formar sobre $\Sigma$, y
$\Sigma^*$ denota la unión de $\Sigma^+$ y $\{\varepsilon\}$.

Dado que los lenguajes son conjuntos, las operaciones sobre conjuntos
también se definen para lenguajes: unión, intersección y complemento
\cite{hopcroft2006introduction}. A continuación se formalizan dos
problemas que se utilizan de manera recurrente en los capítulos
posteriores: determinar si una cadena pertenece a un lenguaje y
determinar si un lenguaje contiene algún elemento.

\begin{definicion}
El \textbf{problema de la palabra} consiste en determinar si una cadena
pertenece a un lenguaje dado.
\end{definicion}

Por ejemplo, dado el lenguaje $L$ de las cadenas binarias que terminan
en $0$, el problema de la palabra para la cadena $1100100$ consiste en
determinar si $1100100 \in L$.

\begin{definicion}
El \textbf{problema del vacío} consiste en determinar si un lenguaje
es vacío.
\end{definicion}

Por ejemplo, el problema del vacío para el lenguaje formado por los
números pares mayores que $5$ que además son primos, consiste en
decidir si tal lenguaje tiene algún elemento. Los antecedentes del
presente trabajo, que se revisan en el capítulo~\ref{cap:rcg-sat},
emplean el problema del vacío para decidir la satisfacibilidad de
fórmulas booleanas, problema que se presenta en la
sección~\ref{sec:sat}.

Una vez introducidos los conceptos básicos y los dos problemas
anteriores, a continuación se presentan las gramáticas, un mecanismo
que permite describir los elementos de un lenguaje mediante reglas de
producción.

\subsection{Gramáticas}
\label{sec:gramaticas}

Los lenguajes formales se pueden describir mediante \textbf{gramáticas},
un mecanismo que permite generar los elementos de un lenguaje a partir
de un símbolo inicial y un conjunto de reglas de producción.

\begin{definicion}
Una \textbf{gramática} es un formalismo utilizado para describir
lenguajes formales. Se define como una $4$-tupla
\[
  G = (N, \Sigma, P, S),
\]
donde
\begin{itemize}
  \item $N$ es un conjunto finito de símbolos llamados
    \textbf{símbolos no terminales};
  \item $\Sigma$ es un conjunto finito de símbolos llamados
    \textbf{símbolos terminales}, que constituyen el alfabeto sobre el
    que se construyen las cadenas del lenguaje. Se debe cumplir que
    $N \cap \Sigma = \emptyset$;
  \item $P$ es un conjunto finito de \textbf{reglas de producción} de la
    forma
    \[
      \alpha \to \beta,
      \quad \text{con } \alpha, \beta \in (N \cup \Sigma)^*;
    \]
  \item $S \in N$ es el \textbf{símbolo inicial} de la gramática.
\end{itemize}
\end{definicion}

Una \textbf{derivación} en la gramática consiste en seleccionar una
regla de producción $\alpha \to \beta$ y sustituir una ocurrencia de
$\alpha$ en una cadena $w \in (\Sigma \cup N)^+$ por $\beta$
\cite{hopcroft2006introduction}.

Por ejemplo, considere la gramática $G = (\{S\}, \{a, b\}, P, S)$ con
producciones $P = \{S \to aSb,\; S \to ab\}$. Partiendo de $S$, una
primera derivación aplica $S \to aSb$ y produce la cadena $aSb$; una
segunda derivación aplica $S \to ab$ sobre la ocurrencia de $S$ en
$aSb$ y produce $aabb$.

Se dice que la gramática $G$ genera una cadena $w \in \Sigma^*$ si
existe una sucesión finita de derivaciones que, partiendo de $S$,
produce $w$. El \textbf{lenguaje generado} por una gramática $G$ se
denota
\[
  L(G) = \{ w \in \Sigma^* \mid S \xrightarrow{+} w \},
\]
donde $\xrightarrow{+}$ indica la existencia de una sucesión finita
de una o más derivaciones desde $S$ hasta $w$. En el ejemplo anterior,
$L(G) = \{a^n b^n \mid n \ge 1\}$.

Una vez presentadas las gramáticas, a continuación se expone la
jerarquía de Chomsky, que clasifica los lenguajes formales según las
restricciones en las reglas de producción de las gramáticas que los
generan.

\subsection{Jerarquía de Chomsky}

La \textbf{Jerarquía de Chomsky}~\cite{hunter2020chomsky} clasifica los
lenguajes en cuatro tipos, según las restricciones en las reglas de
producción de las gramáticas que los generan.

El primer tipo son las \textbf{gramáticas irrestrictas}, que no imponen
restricciones sobre sus reglas de producción. Cada regla tiene la forma
$\alpha \to \beta$, donde $\alpha, \beta \in (N \cup \Sigma)^*$ y
$\alpha \neq \varepsilon$. Todo lenguaje generado por una gramática
irrestricta se denomina \textbf{lenguaje recursivamente
enumerable}~[REF].

El segundo tipo son las \textbf{gramáticas dependientes del contexto}.
En estas, cada regla tiene la forma $\alpha A \gamma \to \alpha \beta \gamma$,
donde $A \in N$, $\alpha, \beta, \gamma \in (N \cup \Sigma)^*$, y
$|\beta| \ge 1$. Todo lenguaje generado por una gramática dependiente
del contexto se denomina \textbf{lenguaje dependiente del contexto}.
Todo lenguaje dependiente del contexto es también un lenguaje
recursivamente enumerable. En~\cite{gomezmolina2025} se demuestra que
el lenguaje de las fórmulas booleanas satisfacibles no es libre del
contexto, por lo que necesariamente se sitúa en alguno de estos dos
tipos superiores; este resultado se discute en el
capítulo~\ref{cap:rcg-sat}.

El tercer tipo son las \textbf{gramáticas libres del contexto} (CFG).
En estas, cada regla tiene la forma $A \to \beta$, donde $A \in N$ y
$\beta \in (N \cup \Sigma)^*$. Todo lenguaje generado por una gramática
libre del contexto se denomina \textbf{lenguaje libre del contexto}.
Todo lenguaje libre del contexto es también un lenguaje dependiente del
contexto~\cite{hopcroft2006introduction}.

Las gramáticas de concatenación de rango, formalismo central del
capítulo~\ref{cap:rcg-sat}, se introducen como una extensión de las
CFG. Además, en~\cite{fernandezarias2007} se emplea una CFG para
resolver instancias del problema de la satisfacibilidad booleana, que
se presenta en la sección~\ref{sec:sat}.

El cuarto tipo son las \textbf{gramáticas regulares}. En estas, las
reglas de producción tienen la forma
\[
  A \to aB \quad \text{o} \quad A \to a,
\]
donde $A, B \in N$ y $a \in \Sigma$.

Todo lenguaje generado por una gramática regular se denomina
\textbf{lenguaje regular}. Todo lenguaje regular es un lenguaje libre
del contexto. Los lenguajes regulares se relacionan con los autómatas
finitos, que se presentan en la sección~\ref{sec:automatas}.

Cada tipo de la jerarquía de Chomsky agrupa los lenguajes generados
por las gramáticas de ese tipo. Estas agrupaciones se denominan
\textbf{clases de lenguajes}: la clase de los lenguajes regulares, la
de los libres del contexto, la de los dependientes del contexto y la
de los recursivamente enumerables. Por brevedad, expresiones como "los
lenguajes libres del contexto" se refieren a la clase correspondiente.

Una clase de lenguajes $\mathcal{A}$ es más expresiva que una clase
$\mathcal{B}$ si todo lenguaje de $\mathcal{B}$ pertenece también a
$\mathcal{A}$ y existe al menos un lenguaje de $\mathcal{A}$ que no
pertenece a $\mathcal{B}$. Por ejemplo, la clase de los lenguajes
libres del contexto es más expresiva que la clase de los lenguajes
regulares~\cite{hopcroft2006introduction}.

La diferencia entre las clases de lenguajes de la jerarquía de Chomsky
se puede ilustrar con una familia\footnote{El término "familia de
lenguajes" se utiliza aquí en un sentido distinto al de "clase de
lenguajes" definido en el párrafo anterior: se refiere a una colección
de lenguajes que se obtiene al variar un parámetro ---en este caso,
los lenguajes $L_{\mathrm{copy}}^k$ al variar $k$--- y no a una clase
de la jerarquía de Chomsky.} de lenguajes derivada del lenguaje Copy.
Dado un alfabeto $\Sigma$, para cada entero $k \ge 1$ se define el
lenguaje
\[
  L_{\mathrm{copy}}^k = \{ w^k \mid w \in \Sigma^* \}.
\]
La familia $\{L_{\mathrm{copy}}^k\}_{k \ge 1}$ es la que ilustra la
jerarquía.
El caso $k = 1$ corresponde a $L_{\mathrm{copy}}^1 = \Sigma^*$, que es
un lenguaje regular. Para $k \ge 2$, el lenguaje $L_{\mathrm{copy}}^k$
no es libre del contexto, aunque sí es dependiente del contexto
\cite{hopcroft2006introduction}. En el capítulo~\ref{cap:rcg-sat} se
emplea este lenguaje para ilustrar el poder expresivo de las gramáticas
de concatenación de rango.

La relación de inclusión entre las clases de la jerarquía se ilustra
en la figura~\ref{fig:chomsky}: una clase está contenida en otra
cuando la elipse que la representa aparece dentro de la otra.

\begin{figure}[H]
\centering
\begin{tikzpicture}[font=\small]
  \node[above, ellipse, minimum height=2em, minimum width=6em,  draw] (a) {Regulares};
  \node[above, ellipse, minimum height=4em, minimum width=12em, draw] (b) {};
  \node[above, ellipse, minimum height=6em, minimum width=18em, draw] (c) {};
  \node[above, ellipse, minimum height=8em, minimum width=24em, draw] (d) {};
  \path (a.north) node[above] {Libres del contexto}
        (b.north) node[above] {Dependientes del contexto}
        (c.north) node[above] {Recursivamente enumerables};
\end{tikzpicture}
\caption{Jerarquía de Chomsky.}
\label{fig:chomsky}
\end{figure}

Cada clase de la jerarquía de Chomsky admite un modelo matemático,
llamado autómata, que reconoce los lenguajes de la clase. En la
sección~\ref{sec:automatas} se define el autómata finito.

\section{Autómatas finitos}
\label{sec:automatas}

Los autómatas correspondientes a las cuatro clases de la jerarquía de
Chomsky son: la máquina de Turing para los lenguajes recursivamente
enumerables, la máquina de Turing linealmente acotada para los
lenguajes dependientes del contexto, el autómata de pila para los
lenguajes libres del contexto y el autómata finito para los lenguajes
regulares~\cite{hopcroft2006introduction}.

\begin{definicion}
Un \textbf{autómata finito}~\cite{hopcroft2006introduction} se define
como una $5$-tupla
\[
  \mathcal{A} = (Q, \Sigma, \delta, q_0, F),
\]
donde $Q$ es un conjunto finito de \textbf{estados}, $\Sigma$ es un
alfabeto llamado \textbf{alfabeto de entrada},
$\delta : Q \times \Sigma \to Q$ es una \textbf{función de transición},
$q_0 \in Q$ es el \textbf{estado inicial} y $F \subseteq Q$ es el
conjunto de \textbf{estados de aceptación}.
\end{definicion}

El autómata comienza en el estado inicial $q_0$ y lee los símbolos de la
cadena uno a uno. En cada paso, la función de transición $\delta$
determina, a partir del símbolo actual y del estado actual, el estado al
que pasa el autómata. Al finalizar la lectura del último símbolo, si el
autómata se encuentra en un estado de $F$, la cadena se acepta; en caso
contrario, se rechaza.

Una vez presentado el autómata finito, en la próxima sección se
introducen los conceptos de complejidad computacional, que se utilizan
en los capítulos posteriores para analizar la complejidad del problema
del vacío en los formalismos estudiados.

\section{Complejidad computacional}
\label{sec:complejidad}

Esta sección se organiza en tres partes. Primero se presenta la
notación asintótica. Después se introduce la reducción polinomial.
Por último se definen las clases de problemas $P$, $NP$ y
$NP$-Completo.

\subsection{Notación asintótica}

La notación asintótica permite describir el crecimiento de una función
$f(n)$ a medida que $n$ tiende a infinito. Se utiliza para clasificar
la dificultad de un problema según el tiempo de ejecución de los
algoritmos que lo resuelven.

\begin{definicion}
Una función $g(n)$ pertenece a $O(f(n))$ si existen constantes positivas
$c$ y $n_0$ tales que
\[
  g(n) \le c \cdot f(n) \quad \text{para todo } n \ge n_0.
\]
\end{definicion}

La definición anterior expresa que $g(n)$ está acotada superiormente
por un múltiplo constante de $f(n)$ para $n$ suficientemente grande.

El tiempo de ejecución de un algoritmo es una función del tamaño $n$
de su entrada. Se dice que un algoritmo tiene \textbf{tiempo
polinomial} si esta función pertenece a $O(n^k)$ para alguna
constante $k$. La \textbf{complejidad computacional} de un problema
se define como el tiempo de ejecución, en el peor caso, del algoritmo
más eficiente que lo resuelve~\cite{hopcroft2006introduction}.

Una vez introducida la noción de tiempo polinomial, en la próxima
subsección se introduce la reducción polinomial como herramienta para
relacionar la dificultad entre problemas.

\subsection{Reducción polinomial}
\label{sec:reduccion}

Los conceptos que se introducen a continuación, y los de la próxima
subsección, se aplican a problemas de decisión. Un \textbf{problema
de decisión} es un problema cuya respuesta para cada instancia es sí
o no.\footnote{Todo problema computacional admite una reformulación
equivalente como problema de decisión~\cite{hopcroft2006introduction};
trabajar con este tipo de problemas simplifica el análisis.}

\begin{definicion}
Una \textbf{reducción polinomial} de un problema de decisión $A$ a un
problema de decisión $B$ es una función $f$ que se calcula en tiempo
polinomial y asigna a cada entrada $x$ de $A$ una entrada $f(x)$ de
$B$, tal que $A$ responde sí sobre $x$ si y solo si $B$ responde sí
sobre $f(x)$~\cite{hopcroft2006introduction}.
\end{definicion}

La existencia de una reducción polinomial de $A$ a $B$ significa que
$B$ es, en términos de complejidad polinomial, al menos tan difícil
como $A$. Como consecuencia, si $B$ se resuelve en tiempo polinomial,
$A$ también. La reducción polinomial es la herramienta con la que se
define la clase $NP$-Completo en la próxima subsección.

\subsection{Clases de problemas}

Los problemas computacionales~\cite{hopcroft2006introduction} se
agrupan en clases según el tiempo de ejecución necesario para
resolverlos o verificar sus soluciones. En este trabajo se emplean las
clases $P$, $NP$ y $NP$-Completo.

\begin{definicion}
Un problema pertenece a la clase $\boldsymbol{P}$ si se puede resolver
en tiempo polinomial~\cite{hopcroft2006introduction}.
\end{definicion}

Por ejemplo, el problema 2-SAT, que se presenta en la
sección~\ref{sec:sat-variantes}, pertenece a $P$.

\begin{definicion}
Un problema pertenece a la clase $\boldsymbol{NP}$ si existe un
algoritmo que, dada una instancia del problema y una solución
candidata, verifica en tiempo polinomial si la solución candidata es
correcta~\cite{hopcroft2006introduction}.
\end{definicion}

Por ejemplo, el problema SAT, que se presenta en la
sección~\ref{sec:sat}, pertenece a $NP$: dada una asignación de
valores a las variables de una fórmula booleana, un algoritmo
polinomial determina si la asignación satisface la fórmula.

\begin{definicion}
Un problema pertenece a la clase $\boldsymbol{NP}$\textbf{-Completo}
si pertenece a $NP$ y, además, cualquier problema en $NP$ se puede
reducir polinomialmente a él~\cite{hopcroft2006introduction}.
\end{definicion}

El SAT del ejemplo anterior es también $NP$-Completo.

La pregunta de si $P = NP$ o $P \neq NP$ permanece sin resolver y
constituye uno de los problemas abiertos centrales de la teoría de la
computación~\cite{hopcroft2006introduction}. El conjunto de problemas
$NP$-Completo proporciona una base para el estudio de esta pregunta:
si algún problema $NP$-Completo se puede resolver en tiempo
polinomial, entonces todos los problemas de $NP$ también.

Una vez introducidas estas clases, en la próxima sección se presenta
el problema de la satisfacibilidad booleana, problema de decisión cuyo
estudio mediante formalismos de la teoría de lenguajes motiva este
trabajo.

\section{Problema de la satisfacibilidad booleana}
\label{sec:sat}

El problema de la satisfacibilidad booleana (SAT) es un problema de
la teoría de la computación y la lógica
matemática~\cite{hopcroft2006introduction}. La sección se organiza
en dos partes. Primero se presentan los elementos del SAT y se da
su definición formal. Después se analiza el estatus del SAT como
problema $NP$-Completo junto con sus variantes polinomiales.

\subsection{Elementos del SAT}
\label{sec:elementos-sat}

A continuación se presentan los elementos del SAT: variables
booleanas, literales, cláusulas y fórmulas en forma normal conjuntiva.

\begin{itemize}
  \item \textbf{Variables booleanas:} una variable booleana es una
    variable que puede tomar uno de dos valores: \textit{verdadero} o
    \textit{falso}.
  \item \textbf{Literales:} un literal es una variable booleana o su
    negación. Formalmente, si $x$ es una variable booleana, entonces
    $x$ y $\neg x$ son literales, donde $\neg x$ representa la
    negación de $x$. Un literal toma los valores \textit{verdadero} o
    \textit{falso} según la asignación de valores a las variables.
  \item \textbf{Cláusulas:} una cláusula es una disyunción, con el
    operador \textsc{or}, de uno o más literales. Por ejemplo, la cláusula
    $(x \vee \neg y \vee z)$ es una disyunción de tres literales: $x$,
    $\neg y$ y $z$. Una cláusula es verdadera si al menos uno de sus
    literales es verdadero; es falsa si todos sus literales son falsos.
  \item \textbf{Fórmulas en forma normal conjuntiva:} una fórmula
    booleana en forma normal conjuntiva, abreviada CNF, es una
    conjunción con el operador \textsc{and} de cláusulas. Por ejemplo,
    \[
      (x \vee \neg y \vee z) \wedge (\neg x \vee y) \wedge
      (x \vee \neg z)
    \]
    es una fórmula en CNF con tres cláusulas y tres variables.
\end{itemize}

Una \textbf{interpretación} de una fórmula booleana es una asignación
de valores \textit{verdadero} o \textit{falso} a cada una de sus
variables. Toda fórmula booleana se puede transformar, en tiempo
polinomial, en una fórmula en CNF que es satisfacible si y solo si la
fórmula original lo es~\cite{hopcroft2006introduction}. Por esta
razón, en el resto del trabajo se asume sin pérdida de generalidad
que toda fórmula booleana está en CNF.

\begin{definicion}
El problema de la \textbf{satisfacibilidad booleana}, o SAT, consiste
en determinar si existe una asignación de valores \textit{verdadero} o
\textit{falso} a las variables de una fórmula booleana, tal que la
fórmula completa sea verdadera.
\end{definicion}

Una fórmula booleana en CNF es \textbf{satisfacible} si existe una
asignación de valores a las variables tal que todas las cláusulas de la
fórmula sean verdaderas simultáneamente.

Una vez definido el problema y sus elementos, a continuación se sitúa
al SAT dentro de la clasificación de complejidad y se describen sus
variantes polinomiales.

\subsection{SAT como problema $NP$-Completo y variantes polinomiales}
\label{sec:sat-variantes}

El SAT es el primer problema demostrado como $NP$-Completo y juega un
rol central en la teoría de la complejidad computacional. El SAT pertenece a la clase
$NP$ porque, dada una asignación de valores a las variables de la
fórmula, se puede verificar en tiempo polinomial si la asignación
satisface la fórmula. La demostración de que el SAT es $NP$-Completo
fue una de las contribuciones principales de Stephen Cook en
1971~\cite{cook1971}, y marcó el inicio de la teoría de la
$NP$-Completitud.

Un SAT en el que cada cláusula tiene exactamente $n$ literales distintos
se denomina $n$-SAT. Para el problema 2-SAT existe una solución
polinomial que determina si la fórmula booleana es satisfacible o no
\cite{halim2013cp3}. Para el problema 3-SAT no se conoce ningún
algoritmo polinomial que permita determinar si una fórmula booleana es
satisfacible o no~\cite{hopcroft2006introduction}. Toda fórmula booleana
del problema $n$-SAT se puede reducir a una fórmula equivalente del
problema 3-SAT, por lo que SAT y 3-SAT son equivalentes en términos de
complejidad computacional~\cite{hopcroft2006introduction}.

Aunque no se conoce algoritmo polinomial para resolver el SAT general,
existen variantes particulares que sí admiten solución en tiempo polinomial,
como 2-SAT, Horn-SAT y XOR-SAT.

El problema \textbf{2-SAT} es una variante del SAT donde cada cláusula
contiene exactamente dos literales. El problema 2-SAT admite solución
en tiempo polinomial mediante una modelación basada en grafos, con
algoritmos como la detección de componentes fuertemente conexas en el
grafo de implicación~\cite{halim2013cp3}.

El problema \textbf{Horn-SAT} es una variante del SAT donde cada
cláusula tiene a lo sumo un literal positivo. El problema Horn-SAT
admite solución en tiempo polinomial mediante el algoritmo de
resolución de Horn~\cite{dowling1984horn}.

El problema \textbf{XOR-SAT} es una variante del SAT donde cada
cláusula representa una operación XOR sobre sus literales. El
problema XOR-SAT admite solución en tiempo polinomial al
transformarlo en un sistema de ecuaciones lineales módulo 2 y aplicar
eliminación de Gauss~\cite{schaefer1978}.

Este trabajo busca ampliar el conjunto de subclases del SAT con
solución en tiempo polinomial mediante las gramáticas de concatenación de rango.
En el próximo capítulo se presentan estas gramáticas, se analizan sus
propiedades y limitaciones, y se revisan los antecedentes directos
que abordan el SAT mediante formalismos de la teoría de lenguajes,
con énfasis en los trabajos previos realizados en la Facultad de
Matemática y Computación de la Universidad de La Habana.

\bibliographystyle{plain}
\bibliography{../bib/referencias}

\end{document}
